\documentclass[10pt,a4paper]{amsart}
\usepackage[margin=19mm]{geometry}
\usepackage{amsmath,amssymb,mathtools}
\usepackage[scr=rsfs,cal=euler]{mathalfa}
\usepackage{xfrac}
\usepackage[unicode,hidelinks,bookmarks=false]{hyperref}
\newcommand{\ie}{i.e.\ }
\newcommand{\cf}{cf.\ }
\newcommand{\resp}{respectively\ }
\newcommand{\Subsection}{\subsection}
\providecommand{\nobreakdash}{\nobreak\hbox{-}}
\setlength{\emergencystretch}{2em}
\multlinegap0pt
\usepackage{mathtools}
\usepackage{bm}
\usepackage{xcolor}
\usepackage{algorithm}\usepackage{algpseudocode}
\usepackage{amssymb}
\usepackage[shortlabels]{enumitem}
\newtheorem{lemma}{Lemma}
\newtheorem{proposition}{Proposition}
\title{Degree bounds and synchronization in Gr\"{o}bner basis computations for affine semi-regular systems}
\author{Momonari Kudo, Kazuhiro Yokoyama}
\date{Source: arXiv:2404.03530v4 (CC BY 4.0)}
\begin{document}
\maketitle

\noindent\emph{Source and licence.} Degree bounds and synchronization in Gr\"{o}bner basis computations for affine semi-regular systems. Version arXiv:2404.03530v4 (CC BY 4.0), distributed by arXiv under the Creative Commons Attribution 4.0 International licence, http://creativecommons.org/licenses/by/4.0/, as recorded in the arXiv licence metadata for this exact version and shown on the abstract page https://arxiv.org/abs/2404.03530v4. Full licence notice: \texttt{licenses/arxiv-2404.03530-CC-BY-4.0.txt}.\par\smallskip

\noindent\textit{Editorial scope.} Counted unit: the article's proposition prop:sync\_partial (source0.tex lines 1308--1333). The article's complete original proof of it is delivered with it (source0.tex lines 1335--1407, one \texttt{proof} environment of 73 lines). No mathematics is written, repaired or re-derived: the statement and the proof are the article's own lines, in the article's own notation, and no retained line is substituted or re-flowed.  Retained uncounted as the source-local context the proof needs: lines 1258--1263.  Nothing was omitted from any retained span, so the package carries no omission record.  No retained line contains a citation, so the wrapper carries no bibliography and no bibliography entry.  No retained line contains a figure environment, an image-inclusion command or drawing code, so no image is delivered and the package declares no asset dependency.  Editorial formatting only, in addition: an AMS \texttt{amsart} preamble that restates the article's own declarations and macros used by the retained lines, namely the environments \texttt{lemma}, \texttt{proposition} on separate counters, together with the standard packages those lines need.  The retained environments print in this document as Lemma 1, Proposition 1 in order of appearance, whereas the article prints the same items with the numbering of its own full document.  The complete native source of the article as distributed in the arXiv e-print for this exact version is bundled unchanged as \texttt{sources/157226/source0.tex}.
\begin{lemma}\label{lm:onestepred}
Let $g,h\in R'$ be homogeneous, and assume that $y\nmid \mathrm{LM}_{\prec^h}(h)$.
If $g^{(a)}$ admits a one-step reduction by $h^{(a)}$, then there is a corresponding one-step reduction of $g$ by $h$. 
Conversely, if $g$ admits a one-step reduction by $h$, then there is a corresponding one-step reduction of $g^{(a)}$ by $h^{(a)}$ for $a\not=0$. 
For $a=0$, if the monomial of the term of $g$ eliminated by $h$ is not divisible by $y$, then there is a corresponding one-step reduction of $g^{(a)}$ by $h^{(a)}$, and otherwise, $g^{(a)}\rightarrow g^{(a)}$ corresponds, which may be called a \textcolor{black}{\em trivial reduction}.
\end{lemma}

\begin{proposition}\label{prop:sync_partial}
Let $a\in K$.
Suppose that, immediately before the next S-pair is selected, the current intermediate bases in the computations for $\langle F^h\rangle$ and $\langle F^h|_{y=a}\rangle$ correspond under $g\mapsto g^{(a)}:=g|_{y=a}$, and that their S-pairs that have neither been processed nor discarded correspond as well.
For $a=0$, the input generators are distinguished by their indices as specified in the computational setting.
Let $\bm{s}=(g_1,g_2)$ be one of these S-pairs in the homogeneous
computation, and put $\bm{s}^{(a)}=(g_1^{(a)},g_2^{(a)})$.
Assume that every polynomial in $\mathcal{G}_{\rm hom}^{(\prec\bm{s})}$ has leading monomial not divisible by $y$.
Then the following hold.
\begin{enumerate}
    \item[(1)] The least common multiples of the leading monomials of $\bm{s}$ and $\bm{s}^{(a)}$ agree.
    Hence, $\bm{s}$ can be selected under the normal selection strategy if and only if $\bm{s}^{(a)}$ can be selected.
    Thus, after either pair is selected in one computation, the other pair can be selected in the other computation.

    \item[(2)] Suppose that a chosen reduction sequence in the homogeneous computation is $S(g_1,g_2)=r_0 \xlongrightarrow[h_0]{}\cdots \xlongrightarrow[h_{k-1}]{}r_k$, where $r_k$ is a remainder by $\mathcal{G}_{\rm hom}^{(\prec\bm{s})}$.
    After specialization, and after omitting trivial steps when $a=0$,
    this gives a reduction sequence of $S(g_1^{(a)},g_2^{(a)})$ whose remainder is $r_k^{(a)}$.

    \item[(3)]
    Conversely, suppose that a chosen reduction sequence of $S(g_1^{(a)},g_2^{(a)})$ in the specialized computation has remainder $r'$.
    This sequence can be lifted to a homogeneous reduction sequence.
    If $a\neq0$, the lifted sequence ends in a homogeneous remainder $r$ satisfying $r^{(a)}=r'$.
    If $a=0$, it reaches a homogeneous polynomial $\widetilde r$ with $\widetilde r^{(0)}=r'$; further reductions of terms whose monomials are divisible by $y$, if necessary, yield a homogeneous remainder $r$ with $r^{(0)}=r'$.
\end{enumerate}
If the homogeneous remainder is zero, or if its specialization is nonzero, the current intermediate bases after the corresponding pairs are processed satisfy $\mathcal{G}_{(y=a)}^{(\preceq\bm{s}^{(a)})}= \mathcal{G}_{\rm hom}^{(\preceq\bm{s})}|_{y=a}$.
If, in addition, the leading monomial of every new nonzero homogeneous remainder is not divisible by $y$, then the S-pairs that have neither been processed nor discarded also correspond after this step.
\end{proposition}

\begin{proof}
For {\rm (1)}, let $US$ be the set of all S-pairs that have neither been processed nor discarded immediately before $\bm{s}$ is selected in the homogeneous computation; in particular, $\bm{s}\in US$.
By the hypothesis on the current computations, $US^{(a)}:=(\bm{u}^{(a)})_{\bm{u}\in US}$ is precisely the family of S-pairs that have neither been processed nor discarded in the specialized computation, with the input indices retained when $a=0$.
For every S-pair $\bm{u}=(h_1,h_2)\in US$, our assumption gives $y\nmid \mathrm{LM}_{\prec^h}(h_i)$ for $i=1,2$.
Thus we have $\operatorname{LT}_{\prec^h}(h_i)^{(a)}=\operatorname{LT}_{\prec}(h_i^{(a)})$ and $\operatorname{LM}_{\prec^h}(h_i)=\operatorname{LM}_{\prec}(h_i^{(a)})$, and hence
\[
\operatorname{lcm} \bigl(\operatorname{LM}_{\prec^h}(h_1),\operatorname{LM}_{\prec^h}(h_2)\bigr) = \operatorname{lcm}\bigl(\operatorname{LM}_{\prec}(h_1^{(a)}), \operatorname{LM}_{\prec}(h_2^{(a)})\bigr).
\]
Suppose first that $\bm{s}=(g_1,g_2)$ is selected in the homogeneous
computation.
As the normal selection strategy is used, for every
$\bm{u}=(h_1,h_2)\in US$ we have
\[
\operatorname{lcm}
\bigl(\operatorname{LM}_{\prec^h}(g_1),
      \operatorname{LM}_{\prec^h}(g_2)\bigr)
\preceq^h
\operatorname{lcm}
\bigl(\operatorname{LM}_{\prec^h}(h_1),
      \operatorname{LM}_{\prec^h}(h_2)\bigr).
\]
Since all these leading monomials and least common multiples belong to
$R$, the restriction of $\prec^h$ to them agrees with $\prec$.
Therefore, the corresponding inequality holds in the specialized
computation, and $\bm{s}^{(a)}$ is among the S-pairs that may be
selected under the normal selection strategy.
The converse follows by the same argument.
Thus, after either $\bm{s}$ or $\bm{s}^{(a)}$ is selected in one
computation, its counterpart may be selected in the other.
Moreover, $S(g_1,g_2)^{(a)}=S(g_1^{(a)},g_2^{(a)})$.

For {\rm (2)}, assume that the following reduction sequence is chosen:
\begin{equation}\label{eq:remainder-seq}
S(g_1,g_2)=r_0 \xlongrightarrow[h_0]{} \cdots \xlongrightarrow[h_{k-1}]{} r_k
\end{equation}
with remainder $r_k$ by $\mathcal{G}^{(\prec \bm{s})}_{\rm hom}$, where each reducer $h_i\in \mathcal{G}^{(\prec \bm{s})}_{\rm hom}$ is homogeneous and satisfies $y\nmid \mathrm{LM}_{\prec^h}(h_i)$.
Applying Lemma~\ref{lm:onestepred} successively to each step in \eqref{eq:remainder-seq}, specialization at $y=a$ yields \begin{equation}\label{eq:remainder-seq-specified}
S(g_1^{(a)},g_2^{(a)})=r_0^{(a)} \xlongrightarrow[h_0^{(a)}]{} \cdots \xlongrightarrow[h_{k-1}^{(a)}]{} r_k^{(a)}.
\end{equation}
For $a\neq0$, \eqref{eq:remainder-seq-specified} is a reduction sequence.
If $a=0$, removing the trivial reductions from \eqref{eq:remainder-seq-specified} gives a reduction sequence whose final polynomial is still $r_k^{(a)}$.

Next we show that $r_k^{(a)}$ is a remainder by $\mathcal{G}^{(\prec \bm{s}^{(a)})}_{(y=a)}$.
If $r_k^{(a)}$ were reducible by some $h'\in\mathcal{G}^{(\prec \bm{s}^{(a)})}_{(y=a)}$, then the assumed correspondence between the current intermediate bases would give an element $h\in\mathcal{G}^{(\prec \bm{s})}_{\rm hom}$ such that $h^{(a)}=h'$.
By Lemma~\ref{lm:onestepred}, $r_k$ would then be reducible by $h$, contrary to the fact that $r_k$ is a remainder.
Thus $r_k^{(a)}$ is a remainder, proving {\rm (2)}.

\medskip

For {\rm (3)}, let
\begin{equation}\label{eq:redseqata}
S(g_1^{(a)},g_2^{(a)})=r_0' \xrightarrow[h_0^{(a)}]{} \cdots \xrightarrow[h_{\ell-1}^{(a)}]{} r_\ell'
\end{equation}
be a chosen reduction sequence in the specialized computation with remainder $r_\ell'$, where each $h_i$ is a corresponding element of $\mathcal{G}^{(\prec \bm{s})}_{\rm hom}$.
Applying Lemma~\ref{lm:onestepred} successively to the one-step reductions in \eqref{eq:redseqata}, we obtain a homogeneous reduction sequence
\begin{equation}\label{eq:redseqhom}
S(g_1,g_2)= \widetilde r_0 \longrightarrow \widetilde r_1 \longrightarrow \cdots \longrightarrow \widetilde r_\ell,
\end{equation}
where $\widetilde r_i^{(a)}=r_i'$ for every $i$.
If $a\neq0$, in the same manner as in (2), it can be shown that $\widetilde r_\ell$ is a remainder.
If $a=0$, the polynomial $\widetilde r_\ell$ need not be a remainder.
Any term of $\widetilde r_\ell$ that is still reducible must, however, have a monomial divisible by $y$; otherwise the corresponding one-step reduction would remain nontrivial after specialization and would contradict the fact that $r_\ell'$ is a remainder.
Reducing only such $y$-divisible terms eventually yields a homogeneous remainder $r$.
Since each of these additional reductions specializes to the identity at $y=0$, we still have $r^{(0)}= \widetilde r_\ell^{(0)}= r_\ell'$.

\medskip

Finally, if the homogeneous remainder is zero, neither computation inserts a new polynomial.
If its specialization is nonzero, the two computations insert corresponding remainders, and hence their current intermediate bases continue to correspond.
If, moreover, every new nonzero homogeneous remainder has leading monomial not divisible by $y$, then specialization preserves its leading monomial and does not annihilate it.
The newly generated S-pairs therefore correspond, and Buchberger's product criterion gives the same decision for them because relative primality of their leading monomials is preserved.
Together with the assumed correspondence for the previously available S-pairs, this proves the last assertion.
\end{proof}

\end{document}
